\chapter{Desenvolvimento}
\label{cap:02}


\section{Introdução}


O projeto surgiu no decorrer de uma aula da disciplina onde a professora orientadora disponibilizou um tempo para que o grupo
analizasse as próprias necessidades e identificasse problemas que poderiam ser solucionados com o desenvolvimento de uma aplicação.
Depois de um debate extenso e algumas anotações, o grupo identificou que um problema comum entre os integrantes era a dificuldade
em encontrar informações sobre hardwares, como processadores, placas de vídeo, memórias, etc. Isso especialmente para quem está 
iniciando no assunto, pois os sites especializados costumam apresentar informações técnicas de forma complexa e de difícil compreensão.

\section{Características da Aplicação}

\subsection{Atores e seus papéis}

\begin{itemize} 
    \item {Administrador (\textit{Admin}):} Possui acesso a mutações e a todo o \textit{backlog} do sistema. Tem privilégios globais para gerenciar o Dicionário Contextual, as entidades de \textit{hardware} (fabricantes, modelos, especificações) e as configurações da plataforma.
    \item {Usuário:} Está limitado ao seu próprio escopo. Possui acesso restrito às funcionalidades voltadas ao cliente final, podendo autenticar-se, gerenciar seu próprio perfil (TDU), realizar comparações, favoritar dispositivos e criar alertas de preço.
\end{itemize}

\subsection{Classes de usuários e características}

\text {Os usuários e suas características são descritas no Quadro~\ref{tab:quadroClasses}.}

\subsection{Classes de usuários e características}

\text {Os usuários e suas características são descritas no Quadro~\ref{tab:quadroClasses}.}

\begin{quadro}[!htbp]
\centering
\caption{Classes de usuários e características}
	\begin{tabular}{ | m{4cm} | m{10cm}| }
		\hline
		\cellcolor{gray!20}\textbf{Classe} & \cellcolor{gray!20}\textbf{Descrição}  \\ \hline
		Administradores   & São os usuários com privilégios elevados, responsáveis pelo gerenciamento do sistema, manutenção do Dicionário Contextual, controle do catálogo de \textit{hardware} (fabricantes, modelos e especificações) e operações de \textit{backoffice}. \\ \hline
		Usuários          & São as pessoas que utilizam a plataforma HardDex no seu dia a dia; aproveitam as funcionalidades de busca semântica, realizam comparações de dispositivos com base no seu Tipo de Uso (TDU) e gerenciam seus alertas de preço. \\ \hline
		Desenvolvedores   & São os profissionais responsáveis pela criação e implementação da plataforma HardDex; desenvolvem o sistema escrevendo o código necessário para garantir que ele seja seguro, escalável e atenda a todos os requisitos funcionais e não funcionais. \\ \hline
	\end{tabular}
	\\ \vspace{0.2cm}
	Fonte: Autoria própria (2026)
	\label{tab:quadroClasses}
\end{quadro}



\section{Arquitetura física}

A arquitetura física adota o modelo cliente-servidor. A visualização das páginas é feita no navegador do usuário, enquanto a execução da aplicação fica por conta da comunicação de um servidor com o banco de dados. A Figura~\ref{fig:arqfisica} apresenta este modelo.

\begin{figure}[!htbp]
	\centering
	\caption{Processo de funcionamento da aplicação.}
    \fbox{\includegraphics[scale=.9]{imagens/fig_arquiteturaFisica.png}}
    \legend{Fonte: Autoria própria (2026)}	
	\label{fig:arqfisica}
\end{figure}


\subsection {Tecnologias utilizadas}

As tecnologias utilizadas foram definidas com base no levantamento de requisitos, no ambiente operacional e nas restrições de projeto e implementação. A seleção visa garantir compatibilidade, escalabilidade, segurança e produtividade no desenvolvimento da aplicação.

\begin{itemize} 
    \item {Plataforma de desenvolvimento:} O ambiente de desenvolvimento integrado (IDE) adotado foi o \textit{ \citeonline{vstudio}}, utilizado por sua compatibilidade com múltiplas plataformas e integração com ferramentas de desenvolvimento.
    \item {\textit{Front-end}:} Trata-se da parte do desenvolvimento responsável pela interface gráfica com o usuário, ou seja, o que ele vê e interage com a aplicação. Foram utilizados:
        \begin{itemize} 
        \item {\textit{\citeonline{vue}:}} Estrutura \textit{\citeonline{java}} acessível, de alto desempenho e versátil para criar interfaces de usuário;
        \item {\textit{\citeonline{nuxt}:}} Estrutura para criar aplicativos que facilita o desenvolvimento, tanto da interface gráfica quanto da lógica e dados, com Vue.js.
        \end{itemize}
    \item {\textit{Back-end}:} Trata-se da parte da aplicação responsável pela lógica, armazenamento e processamento de dados, e a interação com bancos de dados e servidores. Foram utilizados:
        \begin{itemize} 
        \item {\textit{\citeonline{node}:}} Ambiente de execução de multiplataforma, que permite criar servidores, aplicações e programas de automação de tarefas.
        \item {\textit{\citeonline{express}:}} Estrutura para aplicativos em Node.js que fornece um conjunto robusto de recursos para aplicativos \textit{web} e móveis.
        \end{itemize}
    \item {Banco de dados:}
        \begin{itemize} 
        \item {\textit{\citeonline{sqlite}:}} Biblioteca em linguagem C que implementa um mecanismo de banco de dados SQL pequeno, rápido e completo.
        \end{itemize}
    \item {Mapeamento Objeto-Relacional(ORM):}
        \begin{itemize}     
        \item {\textit{\citeonline{sequelize}:}} Mapeamento relacional para diversos bancos de dados (SQLite e outros) que apresenta suporte para carregamento e integração.
        \end{itemize}
\end{itemize}


\section{Arquitetura lógica}

Na arquitetura lógica da aplicação, foram definidos os diferentes componentes e módulos que interagem e se organizam para realizar as funcionalidades da aplicação. A Figura~\ref{fig:arqlogica} apresenta de forma simplificada a estrutura utilizada.

\begin{itemize} 
    \item {\textit{Front-end}:} A camada de interface do usuário foi desenvolvida com o \textit{Vue.js}, utilizando o \textit{Nuxt.js} para estruturação do projeto, o roteamento e a renderização das páginas, tanto no lado do cliente quanto do servidor. 

    \item {\textit{Back-end}:} A lógica foi implementada em \textit{Node.js}, com o uso do \textit{Express.js} para a criação da \textit{API RESTful}.

    \item {Banco de dados:} Os dados da aplicação são armazenados utilizando o \textit{SQLite}.

    \item {Mapeamento Objeto-Relacional(ORM):} No mapeamento da aplicação com o banco de dados, foi utilizado o \textit{Sequelize}. 

    \item {Tecnologia \textit{web} e integração:} A comunicação entre o \textit{front-end} e o \textit{back-end} foi feita por meio de \textit{API RESTful}, utilizando o padrão JSON para a troca de dados. A \textit{API Restful} é um tipo de interface de programação de aplicações (API) que se baseia em recursos (como URLs) e utiliza os métodos (\textit{GET, POST, PUT, DELETE}, etc.) para realizar operações como leitura, criação, atualização e exclusão de dados.
\end{itemize}

\begin{figure}[!htbp]
	\centering
	\caption{Arquitetura lógica com as principais tecnologias.}
	\fbox{\includegraphics[scale=.41]{imagens/fig_arquiteturaLogica.png}}
    \legend{Fonte: Autoria própria (2026)}	
	\label{fig:arqlogica}
\end{figure}

\newpage
\subsection {Diagrama de caso de uso}

Os requisitos foram agrupados, e na fase de análise foram definidos os casos de uso. A Figura~\ref{fig:diagramacasouso} representa o diagrama de caso de uso da aplicação.

\begin{figure}[!htbp]
	\centering
	\caption{Diagrama de caso de uso.}
	\fbox{\includegraphics[scale=.6]{imagens/exemploQuadro.png}}
    \legend{Fonte: Autoria própria (2026)}	
	\label{fig:diagramacasouso}
\end{figure}

\newpage
\subsection {Requisitos funcionais}

Na definição dos requisitos funcionais para o desenvolvimento da aplicação, foram agrupados conforme apresentado no Quadro~\ref{tab:casosreq}. 

\begin{quadro}[!htbp]
\centering
\caption{Casos de Uso x Requisitos funcionais}
    \begin{tabular}{ | m{1.7cm} | m{4cm}| m{9cm} | }
		\hline
		\cellcolor{gray!20}\textbf{No} & \cellcolor{gray!20}\textbf{Caso de Uso / Módulo} & \cellcolor{gray!20}\textbf{Descrição resumida} \\ \hline
		RF-001 & Autenticação Unificada (SSO) & Permitir cadastro e login por meio de credenciais próprias ou OAuth 2.0 (Google, Apple, entre outros), mantendo a sessão sincronizada entre Web e App. \\ \hline
		RF-002 & Gestão de Perfil e TDU & Permitir ao usuário gerenciar seus dados pessoais, histórico de buscas e refazer o quiz de Tipo de Uso (TDU) a qualquer momento. \\ \hline
      	RF-003 & Sincronização Cross-Device & Sincronizar dispositivos favoritados e histórico de comparações em tempo real entre todas as plataformas autenticadas. \\ \hline
      	RF-004 & Quiz de Perfilamento & Executar um fluxo interativo de perguntas (orçamento, finalidade principal e portabilidade) para classificar o usuário em um TDU específico. \\ \hline
      	RF-005 & Busca Semântica & Converter termos de linguagem natural (ex.: "celular com câmera boa e barato") em filtros técnicos estruturados no banco de dados. \\ \hline
      	RF-006 & Tags Baseadas em TDU & Exibir badges visuais nos produtos, priorizando atributos relevantes de acordo com o perfil do usuário (ex.: "Bateria Longa" para "Estudante"). \\ \hline
     	RF-007 & Comparação Multidimensional & Gerar uma matriz comparativa lado a lado de até 4 dispositivos, destacando visualmente os atributos superiores. \\ \hline
      	RF-008 & Dicionário Contextual & Exibir modais e tooltips explicativos para termos técnicos, contendo traduções leigas com apoio visual (ícones e descrições simples). \\ \hline
        RF-009 & Gestão de Entidades & Disponibilizar painel administrativo com operações CRUD para fabricantes, modelos, especificações técnicas e mídias. \\ \hline
        RF-010 & Gestão do Dicionário & Oferecer interface administrativa para criação e atualização das traduções leigas dos termos técnicos. \\ \hline
	\end{tabular}
    \label{tab:casosreq}
    \\ \vspace{0.2cm}
	Fonte: Autoria própria (2026)
\end{quadro}

\subsection {Requisitos de regras de negócio}
Os requisitos de regras de negócios são readequados para a aplicação e apresentados no Quadro~\ref{tab:reqregrasnegocio}.

\begin{quadro}[!htbp]
\centering
\caption{Requisitos de regras de negócio}
	\begin{tabular}{ | m{1.7cm} | m{3cm}| m{10cm} | }
		\hline
		\cellcolor{gray!20}\textbf{No} & \cellcolor{gray!20}\textbf{Área de negócio} & \cellcolor{gray!20}\textbf{Descrição resumida} \\ \hline
        RGN-01 & Perfilamento & O usuário pode refazer o quiz de Tipo de Uso (TDU), mas só pode ter um perfil principal ativo por vez.\\ \hline
        RGN-02 & Catálogo & Um produto só pode ser listado no aplicativo se contiver os dados técnicos essenciais exigidos para a busca semântica.\\ \hline
        RGN-03 & Afiliados & Todo redirecionamento para o varejo deve obrigatoriamente injetar as UTMs e o ID de afiliado da plataforma.\\ \hline
        RGN-04 & Autenticação & Contas originadas via OAuth 2.0 não possuem senha gerenciada no sistema, dependendo do token do provedor.\\ \hline
        RGN-05 & Backoffice & Apenas usuários com privilégios de Administrador podem criar, editar ou excluir termos no Dicionário Contextual.\\ \hline             
	\end{tabular}
	\\ \vspace{0.2cm}
	Fonte: Autoria própria (2026)
	\label{tab:reqregrasnegocio}
\end{quadro}


\subsection {Requisitos não funcionais}
Os requisitos não funcionais são listados no Quadro~\ref{tab:reqnaofunc1}.

\begin{quadro}[!htbp]
\centering
\caption{Requisitos não funcionais}
	\begin{tabular}{ | m{1.7cm} | m{4cm}| m{9cm} | }
		\hline
		\cellcolor{gray!20}\textbf{No} & \cellcolor{gray!20}\textbf{Quanto a(o):} & \cellcolor{gray!20}\textbf{Descrição resumida} \\ \hline
        RNF-01 & Segurança & O processamento de preferências, histórico e dados de sessão deve seguir as diretrizes da LGPD (Lei Geral de Proteção de Dados).\\ \hline
        RNF-02 & Sincronização & O histórico de buscas e dispositivos favoritados deve ser sincronizado em tempo real ou com baixíssima latência entre plataformas.\\ \hline
        RNF-03 & Disponibilidade & Os serviços de catálogo e redirecionamento de compras devem possuir alta disponibilidade (uptime de 99,9\%).\\ \hline
        RNF-04 & Compatibilidade & A plataforma deve ser responsiva na Web e possuir aplicativo perfeitamente funcional em ecossistemas Android e iOS.\\ \hline
        RNF-05 & Usabilidade & O dicionário de hardware e os modais explicativos devem seguir as normas de acessibilidade (WCAG) para contraste e leitores de tela.\\ \hline
	\end{tabular}
	\\ \vspace{0.2cm}
	Fonte: Autoria própria (2026)
	\label{tab:reqnaofunc1}
\end{quadro}

\newpage
\subsection {Especificações funcionais}

Aborda-se agora, o detalhamento das especificações funcionais do sistema na seguinte ordem:

- Especificação UC 001: Acesso/Autenticação.

- Especificação UC 002: Manter usuário(a).

- Especificação UC 003: Manter cliente.

- Especificação UC 004: Manter produto. \\


\newpage
\subsubsection {Especificação UC 0001: Acesso/Autenticação}
\begin{center}
    \begin{tabular}{  | c |  }
        \hline
            \textbf {HARDDEX - COMPARAÇÃO DE HARDWARE}\\ 
            \textbf {ESPECIFICAÇÃO FUNCIONAL: UC 0001 - Acesso/Autenticação}\\ 
        \hline
    \end{tabular}
\end{center}

       \textbf {Descrição Resumida do Caso de Uso} 
       
        Este caso de uso tem por objetivo descrever as funcionalidades pertinentes ao acesso e autenticação dos clientes, onde o usuário pode se cadastrar, fornecendo seu nome de login (e-mail e senha) do usuário, permitindo que suas informações estejam sempre atualizadas.

\begin{quadro}[!htbp]
	\begin{tabular}{  m{4cm}  m{11cm} }
		\cellcolor{gray!10}\textbf {Atores}                & \cellcolor{gray!10} Administrador, Usuário.  \\ \hline
        \textbf {Pré-Condição}          & O Administrador ou Usuário deve ter acesso/autenticação no sistema.  \\ \hline
       
        \cellcolor{gray!10}\textbf {Fluxo Principal}       & \cellcolor{gray!10}\textbf {P01.} O usuário acessa o sistema que inicialmente possui a opção de novo cadastro.
        
        \textbf {P02.} O Sistema apresenta os campos em aberto necessários para realizar o cadastro, tais como nome, e-mail, senha, telefone e outros.
        
        \textbf {P03.} O Sistema verifica se login e senha são validos.
        
        \textbf {P04.} O Sistema exibe a mensagem MSG001.
        
        \textbf {P05.} O Usuário acessa o e-mail enviado acessando o link anexado para confirma o cadastro.
        
        \textbf {P06.} O Usuário é levado a tela inicial do aplicativo para realizar seu primeiro login.
        
        \textbf {P07.} O caso de uso termina. \\ \hline

        \textbf {Fluxo  Alternativos}  & \textbf {A01 [P02] - Dados inválidos}
  
        1. O Sistema exibe a mensagem MSG002.
        
        2. O Sistema solicita que os dados sejam preenchidos nova.
        
        3. Os dados são preenchidos e checados até a confirmação de dados válidos.
        
        4. O caso de uso continua a execução no passo [P03].
        
        \textbf {A02 [P02] - Registro existente}
        
        1. O Sistema exibe a mensagem MSG003.
        
        2. O Sistema confere os dados informados que possuem informações cadastradas e fornece a opção de enviar um e-mail de redefinição de senha.
        
        3. O Sistema envia o e-mail de redefinição de senha, se solicitado.
        
        4. O caso de uso continua a execução no passo [P03].
        
        \textbf {A03 [P02] - E-mail em uso}
        
        1. O Sistema exibe a mensagem MSG004.
        
        2. O Usuário fornece um e-mail diferente.
        
        3. O sistema valida se o e-mail não possui registro.
        
        4. Se a validação for um sucesso, o caso de uso continua no passo [P03].  \\ \hline  
	\end{tabular}
\end{quadro}

\newpage
\begin{quadro}[!htbp]
	\begin{tabular}{  m{4cm}  m{11cm} }  	
        \cellcolor{gray!10}\textbf {Exceções}            & \cellcolor{gray!10}\textbf {E01 [A01.3] - Nome do usuário já registrado}
        
        1. O sistema exibe a mensagem MSG003.
        
        2. Retorna ao passo [P02].
        
        \textbf {E02 [A01.3] - CPF do usuário já registrado}
        
        1. O sistema exibe a mensagem MSG003.
        
        2. Retorna ao passo [P02].
        
        \textbf {E03 [A01.3] - Ausência de dados obrigatórios de usuário}
        
        1. O sistema exibe a mensagem MSG005.
        
        2. Retorna ao passo [A01.2].
        
        \textbf {E04 [P01] - Nome ou CPF de cliente não localizado}
        
        1. O sistema exibe a mensagem MSG005.
        
        2. Retorna ao passo [P01].  \\ \hline
	
        \textbf {Pós-Condição}        & Acesso concedido ou negado.  \\ \hline
  	
        \cellcolor{gray!10}\textbf {Produtos Gerados}    & \cellcolor{gray!10} Não se aplica.  \\ \hline
        
        \textbf {Requisitos}          & REQ-01, REQ-02.  \\ \hline
        
        \cellcolor{gray!10}\textbf {Regras de Negócio}   & \cellcolor{gray!10} RGN-01, RGN-02, RGN-03.  \\ \hline
        
        \textbf {Casos de Uso Incluídos} & Não se aplica.  \\ \hline

        \cellcolor{gray!10} \textbf {Casos de Uso Estendidos} & \cellcolor{gray!10} Não se aplica.  \\ \hline

%        \cellcolor{gray!10} \textbf {Casos de Uso Estendidos} & \cellcolor{gray!10} UC002 - Manter Usuário, UC003 - Manter Cliente, UC004 - Manter Produto, UC007 - Realiza Pedido, UC008 - Manter Promoção.  \\ \hline
        
        \textbf {Listagem de Mensagens}   & MSG001: E-mail de autenticação enviado.
        
        MSG002: Nome ou e-mail inválidos.

        MSG003: Dados informados já possuem cadastro.

        MSG004: E-mail em uso.
          
        MSG005: Ausência de dados obrigatórios de usuário.  \\ \hline
        
        \cellcolor{gray!10} \textbf {Interfaces}   & \cellcolor{gray!10} Interface001: Inicial (Figura~\ref{fig:protInicial}).
        
        Interface002: Login (Figura~\ref{fig:protInicial}).  \\ \hline 
	\end{tabular}
\end{quadro}




\newpage
\subsubsection {Especificação UC 0002: Manter usuário}

\begin{center}
    \begin{tabular}{  | c |  }
        \hline
             \textbf {HARDDEX - COMPARAÇÃO DE HARDWARE}\\
            \textbf {ESPECIFICAÇÃO FUNCIONAL: UC 0002 - Manter Usuário}\\
        \hline
    \end{tabular}
\end{center}

        \textbf {Descrição Resumida do Caso de Uso} 
        
        Este caso de uso descreve as funcionalidades para manter os dados dos usuários no sistema, incluindo cadastro, consulta, atualização e exclusão.
        
\begin{quadro}[!htbp]
	\begin{tabular}{  m{4cm}  m{11cm} }
        \cellcolor{gray!10} \textbf {Atores}                & \cellcolor{gray!10} Administrador, Usuário.  \\ \hline
        \textbf {Pré-Condição}          & O Administrador ou Usuário deve ter acesso/autenticação no sistema.  \\ \hline
       
        \cellcolor{gray!10} \textbf {Fluxo Principal}       & \cellcolor{gray!10} \textbf {P01.} O Administrador ou Usuário informa o cnpj/cpf ou nome de usuário para pesquisar.
        
        \textbf {P02.} O Sistema verifica os dados do usuário que combina com a consulta.
        
        \textbf {P03.} O Administrador ou Usuário seleciona o desejado.
        
        \textbf {P04.} O Sistema apresenta os dados preenchidos à respeito do usuário selecionado.
        
        \textbf {P05.} O Sistema exibe as opções disponíveis que podem:
        
        -para o Administrador: inclusão, atualização, exclusão de um usuário, voltar.
        
        - para o Usuário: atualização, exclusão do usuário, voltar.
        
        \textbf {P06.} O Administrador ou Usuário seleciona a opção voltar.
        
        \textbf {P07.} O caso de uso termina. \\ \hline
	
        \textbf {Fluxo  Alternativos}  & \textbf {A01 [P06] - Inserção de um usuário}
        
        1. O Administrador seleciona a opção de inclusão ou outra opção ou opta por finalizar o caso de uso.
        
        2. O Sistema disponibiliza campos para preenchimento pelo usuário.
        
        3. O Administrador informa os dados: código do usuário, nome do usuário, descrição do login, senha, e-mail.
        
        4. O Sistema registra as informações e apresenta a mensagem MSG001.
        
        5. O caso de uso continua a execução no passo [P04].
        
        \textbf {A02 [P06] - Atualização de um usuário}
        
        1. O Administrador seleciona a opção de atualização ou outra opção ou opta por finalizar o caso de uso.
        
        2. O Sistema habilita para o modo editável as informações pertinentes ao usuário.
        
        3. O Administrador informa as atualizações a serem feitas.
        
        4. O Sistema registra as alterações.
        
        5. O caso de uso continua a execução no passo [P04].  \\ \hline
	\end{tabular}
\end{quadro}

\newpage

\begin{quadro}[!htbp]
	\begin{tabular}{  m{4cm}  m{11cm} }
        \textbf {Fluxo  Alternativos}  & \textbf {A03 [P06] Exclusão de um usuário}
        
        1. O Administrador seleciona a opção de exclusão ou outra opção ou opta por finalizar o caso de uso.
        
        2. O Sistema exibe a mensagem MSG002.
        
        3. O Administrador confirma a exclusão.
        
        4. O Sistema realiza a exclusão e apresenta a mensagem MSG003.
        
        5. O caso de uso continua a execução no passo [P01].  \\ \hline
	
        \cellcolor{gray!10} \textbf {Exceções}            & \cellcolor{gray!10} \textbf {E01 [A01.3] - CPF/CNPJ do usuário já registrado}
        
        1. O sistema exibe a mensagem MSG004.
        
        2. Retorna ao passo [P01].
        
        \textbf {E02 [A01.3] - Nome de usuário já registrado}
        
        1. O sistema exibe a mensagem MSG004.
        
        2. Retorna ao passo [P01].
        
        \textbf {E03 [A01.3] - Ausência de dados obrigatórios de usuário}
        
        1. O sistema exibe a mensagem MSG005.
        
        2. Retorna ao passo [A01.2].
        
        \textbf {E04 [P01] - CPF/CNPJ ou Nome do Usuário não localizado}
        
        1. O sistema exibe a mensagem MSG006.      
        
        2. Retorna ao passo [P01].  \\ \hline
	
        \textbf {Pós-Condição}        & Usuário cadastrado, atualizado ou excluído no sistema.  \\ \hline
  	
        \cellcolor{gray!10} \textbf {Produtos Gerados}    & \cellcolor{gray!10} Não se aplica.  \\ \hline
        
        \textbf {Requisitos}          & REQ-05, REQ-06, REQ-07, REQ-08.  \\ \hline
        
        \cellcolor{gray!10} \textbf {Regras de Negócio}   & \cellcolor{gray!10} RGN-04, RGN-05.  \\ \hline
        
        \textbf {Casos de Uso Incluídos} & Não se aplica.  \\ \hline

        \cellcolor{gray!10} \textbf {Casos de Uso Estendidos} & \cellcolor{gray!10} Não se aplica.  \\ \hline
        
        \textbf {Listagem de Mensagens}   & MSG001: Registro inserido com sucesso.
        
        MSG002: Confirma exclusão?

        MSG003: Registro excluído com sucesso.

        MSG004: Registro existente.

        MSG005: Dados obrigatórios não informados.

        MSG006: Registro não localizado.  \\ \hline

        \cellcolor{gray!10} \textbf {Interfaces}   & \cellcolor{gray!10} Interface001: Cadastro (Figura~\ref{fig:protInicial}).
          \\ \hline        
	\end{tabular}
\end{quadro}




\newpage
\subsubsection{Especificação UC 0003: Manter cliente}
\begin{center}
    \begin{tabular}{  | c |  }
        \hline
            \textbf {HARDDEX - COMPARAÇÃO DE HARDWARE}\\
            \textbf {ESPECIFICAÇÃO FUNCIONAL: UC 0003 - Manter Cliente}\\    
           
        \hline
    \end{tabular}
\end{center}

       \textbf {Descrição Resumida do Caso de Uso}        
        
        Este caso de uso descreve as funcionalidades para manter os dados dos clientes no sistema, incluindo cadastro, consulta, atualização e exclusão.
        
\begin{quadro}[!htbp]
	\begin{tabular}{  m{4cm}  m{11cm} }
        \cellcolor{gray!10} \textbf {Atores}                & \cellcolor{gray!10} Administrador, Usuário.  \\ \hline
        \textbf {Pré-Condição}          & O Administrador ou Usuário deve ter acesso/autenticação no sistema.  \\ \hline
       
        \cellcolor{gray!10} \textbf {Fluxo Principal}       & \cellcolor{gray!10} \textbf {P01.} O Responsável pelo setor informa o cpf/cnpj ou nome do cliente para pesquisar.
        
        \textbf {P02.} O Sistema apresenta os dados dos clientes que combinam com a consulta.
        
        \textbf {P03.} O Responsável pelo setor seleciona o cliente desejado.
        
        \textbf {P04.} O Sistema apresenta os dados preenchidos à respeito do cliente selecionado.
        
        \textbf {P05.} O Sistema exibe as opções disponíveis que podem:
        
        - para o Administrador: inclusão de um cliente, atualização de cliente, exclusão de um cliente, voltar; e,
        
        - para o Usuário: inclusão, atualização de cliente, voltar.
        
        \textbf {P06.} O Administrador ou Usuário seleciona a opção voltar.
        
        \textbf {P07.} O caso de uso termina. \\ \hline
		
        \textbf {Fluxo  Alternativos}  & \textbf {A01 [P06] - Inserção de um cliente}
        
        1. O Administrador ou Usuário seleciona a opção de inclusão ou outra opção ou opta por finalizar o caso de uso.
        
        2. O Sistema apresenta os campos para preenchimento dos dados do cliente.
        
        3. O Administrador ou Usuário informa os dados: CPF/CNPJ, nome, endereço, número, complemento, bairro, cep, cidade, estado, pais, telefone, e-mail, descrição e senha.
        
        4. O Sistema registra as informações e apresenta a mensagem MSG001.
        
        5. O caso de uso continua a execução no passo [P04].
        
        \textbf {A02 [P06] - Atualização de um cliente}
        
        1. O Administrador ou Usuário seleciona a opção de atualização ou outra opção ou opta por finalizar o caso de uso.
        
        2. O Sistema habilita os campos para o modo editável com as informações do cliente selecionado.
        
        3. O Administrador ou Usuário informa as atualizações a serem feitas.
        
        4. O Sistema registra as alterações.
        
        5. O caso de uso continua a execução no passo [P04]. \\ \hline
		\end{tabular}
\end{quadro}

\newpage
\begin{quadro}[!htbp]
	\begin{tabular}{  m{4cm}  m{11cm} }
        \textbf {Fluxo  Alternativos}  & \textbf {A03 [P06] - Exclusão de um cliente}
        
        1. O Administrador ou Usuário seleciona a opção de exclusão ou outra opção ou opta por finalizar o caso de uso.
        
        2. O Sistema exibe a mensagem MSG002.
        
        3. O Administrador confirma a exclusão.
        
        4. O Sistema realiza a exclusão e apresenta a mensagem MSG003.        
        
        5. O caso de uso continua a execução no passo [P01].\\ \hline
        
        \cellcolor{gray!10} \textbf {Exceções}            & \cellcolor{gray!10} \textbf {E01 [A01.3] - CPF/CNPJ do cliente já registrado}
        
        1. O sistema exibe a mensagem MSG004.
        
        2. Retorna ao passo [P01].
        
        \textbf {E02 [A01.3] - Nome do cliente já registrado}
        
        1. O sistema exibe a mensagem MSG004.
        
        2. Retorna ao passo [P01].
        
        \textbf {E03 [A01.3] - Ausência de dados obrigatórios de cliente}
        
        1. O sistema exibe a mensagem MSG005.
        
        2. Retorna ao passo [A01.2].
        
        \textbf {E05 [P01] – CPF/CNPJ ou Nome do cliente não localizado}
        
        1. O sistema exibe a MSG006.
        
        2. Retorna ao passo [P01].  \\ \hline
		
        \textbf {Pós-Condição}        & Cliente cadastrado, atualizado ou excluído no sistema.  \\ \hline
  	
        \cellcolor{gray!10} \textbf {Produtos Gerados}    & \cellcolor{gray!10} Não se aplica.  \\ \hline
        
        \textbf {Requisitos}          & REQ-09, REQ-10, REQ-11, REQ-12.  \\ \hline
        
        \cellcolor{gray!10} \textbf {Regras de Negócio}   & \cellcolor{gray!10} RGN-06.  \\ \hline
        
        \textbf {Casos de Uso Incluídos} & Não se aplica.  \\ \hline
        
        \cellcolor{gray!10} \textbf {Casos de Uso Estendidos} & \cellcolor{gray!10} Não se aplica.  \\ \hline
        
        \textbf {Listagem de Mensagens}   & MSG001: Cliente inserido com sucesso.
        
        MSG002: Confirma exclusão?

        MSG003: Cliente excluído com sucesso.

        MSG004: Cliente existente.
        
        MSG005: Dados obrigatórios não informados.

        MSG006: Cliente não localizado.  \\ \hline

        \cellcolor{gray!10} \textbf {Interfaces}   & Interface001: Cadastro (Figura~\ref{fig:protInicial}).
        \\ \hline        
	\end{tabular}
\end{quadro}



\newpage
\subsection {Diagrama de Atividades}

Para \citeonline{sommerville}, o diagrama de atividades representa a sequência de atividades de um sistema ou processo, utilizando elementos como atividades e ações para ilustrar o fluxo de controle e dados. A Figura~\ref{fig:diagramaAtivAdm} apresenta o diagrama do administrador. O processo do sistema é iniciado com a exibição da tela de login, onde o administrador insere suas credenciais. Após a autenticação, ele pode navegar pelas seções de produtos, consultar usuários e emitir relatórios.


\begin{figure}[!htbp]
	\centering
	\caption{Diagrama de atividades (Administrador).}
	\fbox{\includegraphics[scale=.4]{imagens/exemploQuadro.png}}
    \legend{Fonte: Autoria própria (2026)}
	\label{fig:diagramaAtivAdm}
\end{figure}

\newpage
Já a Figura~\ref{fig:diagramaAtivUsu} ilustra o diagrama de atividades do usuário. O processo do sistema se inicia com a exibição da tela de login, onde o usuário insere suas credenciais. Após a autenticação, ele pode alterar seu cadastro, navegar pela seção de produtos, adicionar itens ao carrinho, realizar a compra e finalizar com o processamento do pagamento.

\begin{figure}[!htbp]
	\centering
	\caption{Diagrama de atividades (Usuário).}
	\fbox{\includegraphics[scale=.4]{imagens/exemploQuadro.png}}
    \legend{Fonte: Autoria própria (2026)}
	\label{fig:diagramaAtivUsu}
\end{figure}


\newpage
\subsection {Diagrama de classes de domínio}

 Conforme \citeonline{sommerville}, o diagrama de classes mostra as classes de objeto no sistema e as associações entre elas, ou seja, representa os objetos e suas interações dentro de um domínio específico. A Figura~\ref{fig:diagramaClasses} mostra o diagrama do sistema, destacando as classes, seu atributos e os relacionamentos entre elas.
 
\begin{figure}[!htbp]
	\centering
	\caption{Diagrama de Classes de Domínio}
	\fbox{\includegraphics[scale=.47]{imagens/exemploQuadro.png}}
    \legend{Fonte: Autoria própria (2026)}
	\label{fig:diagramaClasses}
\end{figure}


\newpage
\section{Protótipos}

Para o desenvolvimento do sistema, foram elaborados diversos protótipos de interface com foco em usabilidade, coerência visual e experiência do usuário.



A Figura~\ref{fig:protInicial} apresenta o protótipo da tela inicial da plataforma.

\begin{figure}[!htbp]
	\centering
	\caption{Protótipo tela inicial.}
	\fbox{\includegraphics[scale=.42]{imagens/exemploQuadro.png}}
    \legend{Fonte: Autoria própria (2026)}	
	\label{fig:protInicial}
\end{figure}
